%=============================================================
%  Plantilla UCUENCA – Configuración de formato
%  Universidad de Cuenca · Facultad de Ingeniería
%  Ing. Henry Maldonado
%  Diseño basado en el Libro de Marca UCUENCA
%=============================================================

%-------------------------------------------------------------
% ENCODING & LANGUAGE
%-------------------------------------------------------------
\usepackage[utf8]{inputenc}
\usepackage[T1]{fontenc}
\usepackage{csquotes}
\usepackage{babel}
\usepackage{comment}
\usepackage{ifthen}
\newboolean{showdetails}
\babelprovide[import,main]{spanish}

%-------------------------------------------------------------
% FONTS
% FiraSans (humanista sans, acorde al Libro de Marca UCUENCA)
%-------------------------------------------------------------
\usepackage[sfdefault,light]{FiraSans}
\usepackage{FiraMono}
\usepackage{microtype}

%-------------------------------------------------------------
% PAGE LAYOUT
%-------------------------------------------------------------
\usepackage[top=2.8cm, bottom=2.8cm, left=3cm, right=2.5cm,
            headheight=16pt]{geometry}

%-------------------------------------------------------------
% GRAPHICS & COLOR
%-------------------------------------------------------------
\usepackage{graphicx}
\usepackage[dvipsnames,table]{xcolor}

%-------------------------------------------------------------
% BRAND COLORS  (Libro de Marca UCUENCA, pp. 26-27)
%-------------------------------------------------------------
\definecolor{UCblue}{RGB}{0,40,86}
\definecolor{UCred}{RGB}{165,16,8}
\definecolor{UCgris}{RGB}{111,111,111}
\definecolor{UCbluelight}{RGB}{230,236,244}
\definecolor{UCbluemid}{RGB}{0,62,126}
\definecolor{codebg}{RGB}{248,248,248}
\definecolor{infoboxbg}{RGB}{232,240,252}
\definecolor{successgreen}{RGB}{39,174,96}
\definecolor{warnyellow}{RGB}{241,196,15}

%-------------------------------------------------------------
% COLORES DE CÓDIGO
%-------------------------------------------------------------
\definecolor{arduino}{HTML}{00A3A9}
\definecolor{astructure}{HTML}{818A42}
\definecolor{avariables}{HTML}{128F8F}
\definecolor{afunctions}{HTML}{DB6B21}
\definecolor{aback}{HTML}{E0E0E2}
\definecolor{amyblue}{rgb}{0.01,0.61,0.98}
\definecolor{amygray}{rgb}{0.47,0.47,0.33}
\definecolor{numb}{rgb}{0.7,0.0,0.7}
\definecolor{punct}{rgb}{0.5,0.5,0.5}
\definecolor{delim}{rgb}{0.3,0.3,0.3}

%-------------------------------------------------------------
% TABLES
%-------------------------------------------------------------
\usepackage{booktabs}
\usepackage{tabularx}
\usepackage{multirow}
\usepackage{array}
\usepackage{longtable}
\usepackage{xltabular}

%-------------------------------------------------------------
% HEADERS & FOOTERS
%-------------------------------------------------------------
\usepackage{fancyhdr}
\usepackage{lastpage}

%-------------------------------------------------------------
% LISTS
%-------------------------------------------------------------
\usepackage{enumitem}
\newcommand{\keys}[1]{%
  {\small\normalfont\sffamily\fboxsep=2pt\fbox{\vphantom{Xy}#1}}}
\newcommand{\ctrl}{Ctrl}
\newcommand{\alt}{Alt}
\newcommand{\shift}{Shift}

%-------------------------------------------------------------
% CODE LISTINGS
%-------------------------------------------------------------
\usepackage{listings}
\usepackage{tcolorbox}
\tcbuselibrary{listings, skins, breakable}

%-------------------------------------------------------------
% MATH
%-------------------------------------------------------------
\usepackage{amsmath,amssymb}

%-------------------------------------------------------------
% HYPERLINKS
%-------------------------------------------------------------
\usepackage[colorlinks=true,
            linkcolor=UCblue,
            urlcolor=UCblue,
            citecolor=UCblue,
            bookmarks=true,
            unicode=true,
            pdftitle={Documento},
            pdfauthor={Ing. Henry Maldonado}]{hyperref}

%-------------------------------------------------------------
% BIBLIOGRAPHY
%-------------------------------------------------------------
\usepackage[backend=biber,style=ieee]{biblatex}
\addbibresource{references.bib}

%-------------------------------------------------------------
% MISC
%-------------------------------------------------------------
\usepackage{caption}
\usepackage{subcaption}
\usepackage{float}
\usepackage{mdframed}
\usepackage{titlesec}
\usepackage{titletoc}
\usepackage{pifont}
\usepackage{fontawesome5}
\usepackage{tikz}
\usetikzlibrary{shapes.geometric,arrows.meta,positioning,shadows,
                fit,backgrounds,calc,shapes,arrows}
\usepackage{comment}
\usepackage[american]{circuitikz}
\usetikzlibrary{arrows.meta, calc, positioning}
\usetikzlibrary{shapes,arrows,positioning}
\usetikzlibrary{positioning}

%=============================================================
%  ESTILO DE LISTADOS DE CÓDIGO
%=============================================================
\newcommand*{\FormatDigit}[1]{\ttfamily\textcolor{black}{#1}}
\lstdefinestyle{FormattedNumber}{%
    literate=*{0}{{\FormatDigit{0}}}{1}%
             {1}{{\FormatDigit{1}}}{1}%
             {2}{{\FormatDigit{2}}}{1}%
             {3}{{\FormatDigit{3}}}{1}%
             {4}{{\FormatDigit{4}}}{1}%
             {5}{{\FormatDigit{5}}}{1}%
             {6}{{\FormatDigit{6}}}{1}%
             {7}{{\FormatDigit{7}}}{1}%
             {8}{{\FormatDigit{8}}}{1}%
             {9}{{\FormatDigit{9}}}{1}%
             {.0}{{\FormatDigit{.0}}}{2}%
             {.1}{{\FormatDigit{.1}}}{2}%
             {.2}{{\FormatDigit{.2}}}{2}%
             {.3}{{\FormatDigit{.3}}}{2}%
             {.4}{{\FormatDigit{.4}}}{2}%
             {.5}{{\FormatDigit{.5}}}{2}%
             {.6}{{\FormatDigit{.6}}}{2}%
             {.7}{{\FormatDigit{.7}}}{2}%
             {.8}{{\FormatDigit{.8}}}{2}%
             {.9}{{\FormatDigit{.9}}}{2}%
             {\ }{{ }}{1},%
}

\lstset{%
  basicstyle=\footnotesize\ttfamily,
  breakatwhitespace=false,
  breaklines=true,
  captionpos=b,
  commentstyle=\color{UCgris},
  deletekeywords={...},
  escapeinside={\%*}{*)},
  extendedchars=true,
  keepspaces=true,
  keywordstyle=[1]\color{astructure},
  keywordstyle=[2]\color{avariables},
  keywordstyle=[3]\color{afunctions},
  keywordstyle=[4]\bfseries\color{afunctions},
  language=c++,
  morekeywords={*,...},
  numbers=left,
  numbersep=5pt,
  numberstyle=\tiny\color{amygray},
  rulecolor=\color{black},
  rulesepcolor=\color{amyblue},
  showspaces=false,
  showstringspaces=false,
  showtabs=false,
  tabsize=2,
  title=\lstname,
  emphstyle=\color{avariables},
  frame=single,
  framexleftmargin=15pt,
  rulecolor=\color{arduino},
  literate=%
    {á}{{\'{a}}}1 {é}{{\'{e}}}1 {í}{{\'{i}}}1 {ó}{{\'{o}}}1 {ú}{{\'{u}}}1
    {Á}{{\'{A}}}1 {É}{{\'{E}}}1 {Í}{{\'{I}}}1 {Ó}{{\'{O}}}1 {Ú}{{\'{U}}}1
    {ñ}{{\~{n}}}1 {Ñ}{{\~{N}}}1
}

\lstdefinestyle{Arduino}{%
    style=FormattedNumber,
    keywords={setup, loop, if, else, for, switch, while, do, break, continue, return, goto},
    morekeywords=[2]{HIGH, LOW, INPUT, OUTPUT, INPUT_PULLUP, LED_BUILTIN, true, false,
                     int, float, void, boolean, char, word, long, short, double, string, array},
    morekeywords=[3]{const, pinMode, digitalWrite, digitalRead, analogReference, analogRead,
                     analogWrite, analogReadResolution, analogWriteResolution, tone, noTone,
                     shiftOut, shiftIn, pulseIn, millis, micros, delay, delayMicroseconds,
                     min, max, abs, constrain, map, pow, sqrt, sin, cos, tan, randomSeed,
                     random, lowByte, highByte, bitRead, bitWrite, bitSet, bitClear, bit,
                     attachInterrupt, detachInterrupt, interrupts, noInterrupts,
                     Stream, Keyboard, Mouse, begin, println, print},
    morekeywords=[4]{Serial},
    morecomment=[l]{//},
    morecomment=[s]{/*}{*/},
    emph={const},
}

% Entorno para bloques de código Arduino
\newtcblisting{ArduinoSketchBox}[2][]{
  colframe=arduino,
  enhanced, drop shadow, sharp corners, boxrule=1pt,
  width=\dimexpr\linewidth+\parindent\relax,
  left skip=-\parindent,
  arc=3pt, outer arc=3pt,
  listing only,
  listing options={
    frame=,
    style=Arduino,
  },
  title=#2,
  #1
}

% Enlace externo (ícono)
\newcommand{\ExternalLink}{%
    \tikz[x=1.2ex, y=1.2ex, baseline=-0.05ex]{%
        \begin{scope}[x=1ex, y=1ex]
            \clip (-0.1,-0.1)
                --++ (-0, 1.2)
                --++ (0.6, 0)
                --++ (0, -0.6)
                --++ (0.6, 0)
                --++ (0, -1);
            \path[draw,
                line width = 0.5,
                rounded corners=0.5]
                (0,0) rectangle (1,1);
        \end{scope}
        \path[draw, line width = 0.5] (0.5, 0.5) -- (1, 1);
        \path[draw, line width = 0.5] (0.6, 1) -- (1, 1) -- (1, 0.6);
        }
    }

\newcommand{\doclink}[2]{\href{#1}{#2}\footnote{\url{#1}}}

%=============================================================
%  TCOLORBOX STYLES (plantilla UCUENCA)
%=============================================================
\newtcolorbox{objetivebox}{
  enhanced, breakable,
  colback=UCbluelight,
  colframe=UCblue,
  fonttitle=\bfseries\small\color{white},
  title={\faBookOpen\hspace{5pt}Objetivos de Aprendizaje},
  arc=2pt,
  left=8pt, right=8pt, top=6pt, bottom=6pt,
  boxrule=0pt, titlerule=0pt,
  attach boxed title to top left={yshift=-2pt,xshift=0pt},
  boxed title style={colback=UCblue, arc=2pt, boxrule=0pt},
}

\newtcolorbox{materialbox}{
  enhanced, breakable,
  colback=white,
  colframe=UCred,
  fonttitle=\bfseries\small\color{white},
  title={\faTools\hspace{5pt}Materiales y Equipos},
  arc=2pt,
  left=8pt, right=8pt, top=6pt, bottom=6pt,
  boxrule=0pt,
  attach boxed title to top left={yshift=-2pt},
  boxed title style={colback=UCred, arc=2pt, boxrule=0pt},
}

\newtcolorbox{infobox}[1][Nota]{
  enhanced, breakable,
  colback=infoboxbg,
  colframe=UCblue,
  fonttitle=\bfseries\small,
  title={\faInfoCircle\hspace{5pt}#1},
  arc=2pt, boxrule=0.6pt,
  left=8pt, right=8pt, top=4pt, bottom=4pt,
}

\newtcolorbox{warningbox}[1][Advertencia]{
  enhanced, breakable,
  colback=warnyellow!15,
  colframe=warnyellow!70!black,
  fonttitle=\bfseries\small,
  title={\faExclamationTriangle\hspace{5pt}#1},
  arc=2pt, boxrule=0.6pt,
  left=8pt, right=8pt, top=4pt, bottom=4pt,
}

\newtcolorbox{entregabox}{
  enhanced, breakable,
  colback=successgreen!8,
  colframe=successgreen!70!black,
  fonttitle=\bfseries\small\color{white},
  title={\faClipboardCheck\hspace{5pt}Entregables},
  arc=2pt, boxrule=0pt,
  attach boxed title to top left={yshift=-2pt},
  boxed title style={colback=successgreen!70!black, arc=2pt, boxrule=0pt},
  left=8pt, right=8pt, top=6pt, bottom=6pt,
}

\newtcolorbox{exercisebox}{
	enhanced, breakable,
	colback=UCblue!8,
	colframe=UCblue,
	fonttitle=\bfseries\small\color{white},
	title={\faDumbbell\hspace{5pt}Ejercicio},
	arc=2pt,
	left=8pt, right=8pt, top=6pt, bottom=6pt,
	boxrule=0pt,
	attach boxed title to top left={yshift=-2pt},
	boxed title style={colback=UCblue, arc=2pt, boxrule=0pt},
}

%=============================================================
%  CHAPTER & SECTION FORMATTING  (estilo Libro de Marca)
%=============================================================
\titleformat{\chapter}[display]
  {\normalfont\LARGE\bfseries\color{UCblue}}
  {\color{UCred}\large\bfseries\chaptername\ \thechapter}
  {4pt}
  {}
  [{\color{UCblue}\vspace{-4pt}\rule{\linewidth}{2pt}}]
\titlespacing*{\chapter}{0pt}{-20pt}{28pt}

\titleformat{\section}
  {\normalfont\large\bfseries\color{UCblue}}
  {\thesection}{1em}{}
  [{\color{UCred}\vspace{1pt}\rule{\linewidth}{1.5pt}}]

\titleformat{\subsection}
  {\normalfont\normalsize\bfseries\color{UCred}}
  {\thesubsection}{1em}{}

\titleformat{\subsubsection}
  {\normalfont\normalsize\bfseries\color{UCblue}}
  {\thesubsubsection}{1em}{}

%=============================================================
%  HEADER / FOOTER  (estilo Libro de Marca UCUENCA)
%=============================================================
\pagestyle{fancy}
\fancyhf{}
\fancyhead[L]{%
  \footnotesize\bfseries\color{UCred}\MakeUppercase{Título del Documento}}
\fancyhead[R]{%
  \footnotesize\bfseries\color{UCblue}\MakeUppercase{Universidad de Cuenca}}
\fancyfoot[R]{\small\color{UCgris}\thepage}
\renewcommand{\headrulewidth}{0.4pt}
\renewcommand{\footrulewidth}{0pt}

\fancypagestyle{plain}{
  \fancyhf{}
  \fancyfoot[R]{\small\color{UCgris}\thepage}
  \renewcommand{\headrulewidth}{0pt}
}

%=============================================================
%  MACROS DE UTILIDAD
%=============================================================
\newcommand{\pinref}[1]{\texttt{#1}}
\newcommand{\comp}[1]{\textbf{#1}}
\newcommand{\career}{Ingeniería en Telecomunicaciones}

% Tikz settings para diagramas de control
\tikzset{
  block/.style = {draw, fill=white, rectangle, minimum height=3em, minimum width=3em},
  tmp/.style  = {coordinate},
  sum/.style  = {draw, fill=white, circle, node distance=1cm},
  input/.style = {coordinate},
  output/.style= {coordinate},
  pinstyle/.style = {pin edge={to-,thin,black}}
}
